\section{The count of elliptic pairs}\label{sec:elliptic-counting}

We now count the elliptic pairs of Section~\ref{sec:setting} on a fixed
lattice \(\Lambda\subset\CC\), and put \(g_i=g_i(\Lambda)\). The pole is
at zero, and the constant term of the profile is retained.
Neither reflection nor a lattice automorphism is
factored out: pairs related in this way are counted separately (so that
\(M_3=8\) below, and not fewer). Two lattices differ by a rescaling
exactly when they have the same invariant
\(j(\Lambda)=1728g_2^3/(g_2^3-27g_3^2)\).

We write elliptic functions through the Weierstrass relations
\begin{equation}\label{eq:weierstrass}
 Y^2=4X^3-g_2X-g_3,\qquad
 \mathscr D X=Y,\qquad
 \mathscr D Y=6X^2-\frac{g_2}{2}.
\end{equation}
For \(g_2^3\ne27g_3^2\), that is, on a lattice, these relations are
realized by \(X=\wp(z;g_2,g_3)\), \(Y=\wp'(z;g_2,g_3)\) and
\(\mathscr D=\mathrm d/\mathrm dz\); see \citet[Chapter~23]{DLMF}.
We use \(X\) and \(Y\) because the formulas below must keep their
meaning for the degenerate parameters \(g_2^3=27g_3^2\), which belong to
no lattice. Using \eqref{eq:weierstrass}, every polynomial in \(X\) and
\(Y\) reduces to a unique normal form \(F(X)+YG(X)\). We say that
\(X^kY^\epsilon\), \(\epsilon\in\{0,1\}\), has pole order
\(2k+3\epsilon\); on a lattice this is its pole order at zero.

The representation of the profile is the classical one. An elliptic
function with one pole per period has zero residue there.
The functions \(1,\wp,\wp',\ldots,\wp^{(p-2)}\) have distinct
principal parts and span the space of elliptic functions whose only
pole is at zero, of order at most \(p\).
Subtracting the indicated principal part from a one-pole solution
leaves an entire elliptic function, hence a constant.
Since \(D^{p-2}\wp=(-1)^{p}(p-1)!\,z^{-p}+O(1)\), the leading
coefficient \eqref{eq:leadingpole} shows that every profile has the form
\begin{equation}\label{eq:elliptic-enumeration-profile}
 V=-s_p\mathscr D^{p-2}X
       +\sum_{j=0}^{p-3}u_j\mathscr D^jX+b.
\end{equation}
Its only pole on \(\CC/\Lambda\) is zero, of order \(p\).
In particular \(\Lambda\) is its full period lattice.
The matching equations \(\mathcal E_p(g_2,g_3)\) are obtained by formal
substitution of \eqref{eq:elliptic-enumeration-profile} in
\eqref{eq:main}, using only \eqref{eq:weierstrass}: the left side is
reduced to its normal form, and the coefficients of \(F\) and \(G\) are
set equal to zero. The unknowns are \(a_0,\ldots,a_{p-1}\),
\(u_0,\ldots,u_{p-3}\) and \(b\). The equations are polynomials in the
unknowns and in \(g_2,g_3\), and they are defined for every
\((g_2,g_3)\in\CC^2\), the degenerate parameters included; on a lattice
their solutions are the elliptic pairs. The symbol
\(\mathcal E_p(g_2,g_3)\) also denotes the set of solutions.
Multiplicity and simplicity are understood as
in Section~\ref{sec:counting}.

\begin{theorem}[Fixed-lattice elliptic enumeration]\label{thm:elliptic-count}
Let \(2p-1\) be prime.
\begin{enumerate}
\item[\textup{(a)}] On every lattice the elliptic pairs are
finite in number, and counted with multiplicity there are
\begin{equation}\label{eq:elliptic-length}
 M_p:=2\binom{2p-2}{p-2}=2S_p
\end{equation}
of them.
\item[\textup{(b)}] There is a proper algebraic subset \(\mathcal N\) of
the \((g_2,g_3)\)-plane outside which all of them are simple. The
lattices in \(\mathcal N\) are those whose invariant \(j(\Lambda)\) lies
in a finite set, which depends on \(p\). Thus a generic lattice, that
is, a lattice whose invariant avoids these values, admits exactly
\(M_p\) distinct elliptic pairs.
\end{enumerate}
\end{theorem}

\paragraph{The smallest case.}
For \(p=2\) the profile is \(V=-12X+b\), and the left side of
\eqref{eq:main} reduces to
\[
 -12a_1Y-12(a_0+b)X+6g_2+a_0b+\frac12b^2.
\]
Consequently the complete elliptic system is
\begin{equation}\label{eq:kdvb-elliptic}
 a_1=0,\qquad
 v=-12\wp-a_0,\qquad a_0^2=12g_2.
\end{equation}
There is no one-pole elliptic member with nonzero Burgers coefficient:
both elliptic pairs belong to the KdV equation, in agreement with
\(M_2=2\). They are the two backgrounds \eqref{eq:shift} of one
cnoidal wave. For \(g_2\ne0\) they are distinct and simple.
For \(g_2=0\), \(g_3\ne0\), they coincide, and the lattice carries one
pair of multiplicity two. Here \(\mathcal N\) is the line \(g_2=0\),
and the finite set of (b) is \(\{0\}\). At order three it is
\(\{0,-300\}\) (Section~\ref{sec:ks}).

\paragraph{No escape to infinity.}
The number \(M_p\) is twice the number \(S_p\) of pulses of
Proposition~\ref{cor:endpoints}, and the proof shows why.
When the parameters of a square polynomial system (as many equations as
unknowns) vary, its solutions may merge or separate. As long as they
remain finite in number, their number, counted with multiplicity, can
change only if a solution moves off to infinity. A solution that moves
off to infinity, rescaled to size one, converges to a nonzero solution
of the leading parts of the equations; this is the rescaling argument
of Step~3 in the proof of Theorem~\ref{thm:prime}, with the ordinary
absolute value in place of the \(\ell\)-adic one. Here the leading
parts are the equations with both periods infinite
(\(\wp\to z^{-2}\); the cusp \(g_2=g_3=0\)), and by rigidity that
rational problem has only the trivial solution. So nothing escapes, and
the number of pairs, counted with multiplicity, is the same for all
\((g_2,g_3)\). We find it with one period infinite (\(\wp\) becomes a
hyperbolic function; the node): there every elliptic pair becomes a
pulse on one of its two constant backgrounds, and every pulse on either
background arises in this way exactly once. The pulses are simple, and
simplicity spreads from the node to a generic lattice.

\begin{lemma}[Conservation of number]\label{lem:conservation}
Let \(E_1,\ldots,E_n\) be polynomials in unknowns
\(x=(x_1,\ldots,x_n)\) and parameters \(t=(t_1,\ldots,t_m)\),
weighted-homogeneous for some positive integer weights of the \(x_j\)
and the \(t_k\). Suppose that the leading system \(E_i(x;0)=0\),
\(1\le i\le n\), has no solution in \(\CC^n\) other than \(x=0\). Then:
\begin{enumerate}
\item[\textup{(i)}] for every \(t\) the system \(E(x;t)=0\) has
finitely many solutions, and they stay bounded while \(t\) stays
bounded;
\item[\textup{(ii)}] the number of solutions, counted with
multiplicity, is the same for every \(t\in\CC^m\);
\item[\textup{(iii)}] the \(t\) for which some solution is not simple
form an algebraic subset of \(\CC^m\), that is, the common zero set of
finitely many polynomials in \(t\).
\end{enumerate}
\end{lemma}

The proof of the lemma uses more algebra than the rest of the paper and
can be skipped; Steps~1--5 below use only its statement.

\begin{proof}
Since the leading system has only the zero solution, the polynomials
in \(x\) modulo its equations form a finite-dimensional space
\citep[Chapter~5, \S3]{CoxLittleOSheaIVA}, spanned by monomials
\(m_1=1,\ldots,m_r\). Every polynomial \(f\) in \((x,t)\) is then,
modulo the equations \(E_i\), a combination \(\sum_ic_i(t)m_i\) with
polynomials \(c_i\) in \(t\). This is proved by induction on the
weight of \(f\), which may be taken weighted-homogeneous: write
\(f=f(x,0)+\sum_kt_kf_k\), write \(f(x,0)\) as a combination of the
\(m_i\) plus \(\sum_ih_iE_i(x;0)\) with weighted-homogeneous \(h_i\),
and replace each \(E_i(x;0)\) by \(E_i\) minus its terms that contain
some \(t_k\). What remains is a sum of products of a \(t_k\) with a
polynomial of lower weight, because \(t_k\) has positive weight.
In particular \(x_jm_i=\sum_lM^{(j)}_{il}(t)m_l\) modulo the equations.
Multiplying this vector identity by the adjugate of the matrix
\(x_j-M^{(j)}(t)\) shows that \(\det\bigl(x_j-M^{(j)}(t)\bigr)\,m_i\) is
a combination of the \(E_i\) for every \(i\) (the Cayley--Hamilton
argument); for \(m_1=1\) this is the determinant itself.
So at every solution \(x_j\) is a root of a
monic polynomial of degree \(r\) whose coefficients are polynomials in
\(t\). This proves~(i).

For (ii), join two parameter values by a segment. By (i) some ball
\(|x|<\varrho\) contains all solutions for all \(t\) on the segment, so
no solution crosses the sphere \(|x|=\varrho\). The number of zeros of
a holomorphic map in a ball, counted with multiplicity, when they are
finite in number and none lies on the bounding sphere, is the degree
of \(E/|E|\) on that sphere, and it does not change under a
deformation without zeros on the sphere
\citep[Appendix~B]{Milnor1968},
\citep[Sections~5.2 and~5.4]{ArnoldGuseinZadeVarchenko1985}; this is
Rouch\'e's theorem in several variables.

For (iii), the set in question is the projection to \(\CC^m\) of the
algebraic set on which \(E_1,\ldots,E_n\) and the Jacobian determinant
\(\det(\partial E_i/\partial x_j)\) vanish.
The monic polynomials of (i) are combinations of the \(E_i\). Eliminate
the unknowns one at a time: at each step the polynomial of (i) for the
unknown being eliminated has leading coefficient \(1\), so by the
Extension Theorem of elimination theory
\citep[Chapter~3, \S1]{CoxLittleOSheaIVA} every solution of the
eliminated equations extends. The projection is therefore the common
zero set of the polynomials in \(t\) alone that are combinations of the
\(E_i\) and the Jacobian determinant.
\end{proof}

Three operations do not change multiplicities, and the proof of the
theorem uses them freely: an invertible affine change of the unknowns;
invertible linear combinations of the equations with constant
coefficients; and the elimination of an unknown that one of the
equations gives as a polynomial in the others. Each of them leaves
unchanged, up to an isomorphism, the space of power series modulo the
combinations of the equations, whose dimension is the multiplicity
(Section~\ref{sec:counting}).

\begin{proof}[Proof of Theorem~\ref{thm:elliptic-count}]
\emph{Step 1: eliminating the operator.}
The coefficient at pole order \(2p\) vanishes by the choice of
\(s_p\).
The coefficients at orders \(2p-1,\ldots,p\) determine
\(a_{p-1},\ldots,a_0\) successively.
At each step the coefficient of the new \(a_j\) is the nonzero
rational leading coefficient of \(D^jV\).
This is a polynomial elimination: each \(a_j\) is obtained as a
polynomial in the remaining unknowns and in \(g_2,g_3\), and
multiplicities do not change.
The residual has pole order at most \(p-1\) and zero residue, since
no normal form has pole order one. It lies in the span of \(1\) and the
\(\mathscr D^jX\), \(0\le j\le p-3\).
Hence \(p-1\) equations \(E_1=\cdots=E_{p-1}=0\) remain in the \(p-1\)
profile coefficients \(u_j,b\): the coefficients of the residual in this
basis.

\emph{Step 2: weights.}
Give \(D,X,g_2,g_3\) weights \(1,2,4,6\), respectively, and give
the profile weight \(p\); this is the usual homogeneity of \(\wp\)
under a rescaling of \(z\).
Then \(u_j\) has weight \(p-j-2\), and \(b\) has weight \(p\).
After the triangular elimination, the residual principal-part
coefficients and the residual constant, that is, the \(E_i\), are
weighted-homogeneous polynomials in the \(u_j\), \(b\), \(g_2\),
\(g_3\), of weights
\begin{equation}\label{eq:elliptic-weights}
 \begin{array}{c|c}
 \text{variables}&1,2,\ldots,p-2,p\\
 \text{equations}&p+1,p+2,\ldots,2p-2,2p.
 \end{array}
\end{equation}
The B\'ezout number of this table is
\((p+1)\cdots(2p-2)(2p)/\bigl((p-2)!\,p\bigr)=2\binom{2p-2}{p-2}=M_p\)
(at \(p=3\): \(4\cdot6/(1\cdot3)=8\)), in agreement with the count
below (companion note).
The normal forms and the eliminations are polynomial in
\(g_2,g_3\), so these statements extend to the degenerate parameters.
For fixed \(g_2,g_3\), the highest weighted parts of the equations in
the profile variables are obtained by setting \(g_2=g_3=0\), since
\(g_2\) and \(g_3\) carry positive weight.

\emph{Step 3: the cusp.}
At the cusp \(g_2=g_3=0\) the relations \eqref{eq:weierstrass} are
satisfied by \(X=z^{-2}\), \(Y=-2z^{-3}\) and
\(\mathscr D=\mathrm d/\mathrm dz\), so a solution of the cusp system
gives a nonconstant rational solution \(V\) of \eqref{eq:main} whose
only pole is at zero. By Lemma~\ref{lem:laplace}, \(V\) vanishes at
infinity, so that
\(b=0\); moreover \(a_0=P(0)=0\), and \(V=(-1)^{p-1}s_pA^\sharp\) for a
monic \(A\) of degree \(p-1\) that divides \(H_A\).
By rigidity, which holds by Theorem~\ref{thm:prime}(c), \(A=\rho^{p-1}\).
Thus \(V=c_\ast z^{-p}\), and all the variables in
\eqref{eq:elliptic-enumeration-profile} vanish, because the functions
\(D^jz^{-2}\) have different pole orders.
The leading equations have no nonzero common zero.

\emph{Step 4: conservation of number.}
The residual equations \(E_1,\ldots,E_{p-1}\) of Step~1
are polynomials in the \(p-1\) profile unknowns and in \(g_2,g_3\),
weighted-homogeneous for the weights of Step~2, and by Step~3 their
leading system, which is the cusp system, has only the zero solution.
Lemma~\ref{lem:conservation} applies, with \(t=(g_2,g_3)\). By (i),
every system \(\mathcal E_p(g_2,g_3)\) has finitely many solutions,
which proves the finiteness in (a). By (ii), their number, counted with
multiplicity, is the same at every point of
the \((g_2,g_3)\)-plane, the degenerate parameters included; at the
cusp all of this multiplicity sits
at the single solution where every unknown vanishes. By (iii), the
parameters at which some solution is not simple form an algebraic
subset \(\mathcal N\) of the plane.

\emph{Step 5: the node.}
To find the common number, and to show that \(\mathcal N\) is not the
whole plane, consider the nodal parameters
\begin{equation}\label{eq:nodal-parameters}
 g_2=\frac1{12},\qquad g_3=-\frac1{216}.
\end{equation}
For these parameters the relations \eqref{eq:weierstrass} are
satisfied by
\[
 X=Q^2-Q+\frac1{12},\qquad
 Y=(2Q-1)Q(1-Q),\qquad \mathscr D=D\ \text{ with }\ DQ=Q(1-Q);
\]
this is the degenerate \(\wp\)-function with the single period
\(2\pi i\), with its pole moved to \(z=i\pi\), written in the logistic
coordinate. A monomial \(X^kY^\epsilon\) becomes a polynomial in
\(Q\) of degree \(2k+3\epsilon\), which is also its pole order. These
degrees are all different, so the monomials remain linearly
independent, and the matching equations for the nodal parameters are
exactly the conditions for \eqref{eq:main} to hold as an identity
between polynomials in \(Q\).
Since \(X=\frac1{12}-DQ\) and \(\mathscr D^jX=-D^{j+1}Q\) for
\(j\ge1\), the profile \eqref{eq:elliptic-enumeration-profile}
becomes, with \(u_{p-2}=-s_p\),
\[
 V=\beta+v_A,\qquad
 A(\rho)=-\frac1{s_p}\sum_{j=0}^{p-2}u_j\rho^{j+1},\qquad
 \beta=b+\frac{u_0}{12}.
\]
Here \(A\) is monic of degree \(d=p-1\) with \(A(0)=0\), and \(\beta\)
is the common limit of \(V\) at the two ends. With the coefficients
\(\alpha_j\) of \(A\) as in \eqref{eq:normalized-A}, the passage from
\((u_0,\ldots,u_{p-3},b)\) to
\((\alpha_1,\ldots,\alpha_{d-1},\beta)\) is an invertible affine
change of unknowns with rational coefficients.

In these unknowns the constant term of \(P(D)V+V^2/2\) is
\(P(0)\beta+\beta^2/2\), and the remaining terms are
\((P+\beta)(D)v_A+v_A^2/2\). By the proof of
Proposition~\ref{prop:divisibility} the latter equals
\((s_p/2)\bigl[2(P+\beta)A-s_pS_A\bigr](D)Q\), and the elimination of
Step~1 is the division of \(S_A\) by \(A\): it
gives \(P+\beta=s_pC_A/2\), a polynomial in the unknowns, which we
still denote by \(P_A\), and it leaves \(-(s_p^2/2)R_A(D)Q\).
Since \(A(0)=0\), the coefficient \(R_0\) vanishes identically, by
\eqref{eq:constant-remainder}. With \(P(0)=P_A(0)-\beta\), the nodal
matching equations are therefore invertible linear combinations, with
constant coefficients, of
\begin{equation}\label{eq:nodal-matching}
 R_1=\cdots=R_{d-1}=0,\qquad
 \beta\bigl(\beta-2P_A(0)\bigr)=0,
\end{equation}
together with \(P=P_A-\beta\). These are the pulse equations
\eqref{eq:pulse-equations} and one quadratic equation for the
background, and the two systems have the same solutions with the same
multiplicities.

By Proposition~\ref{cor:endpoints} the pulse equations have exactly
\(S_p\) solutions, all simple, with \(P_A(0)\ne0\) at each of them.
For each of them the quadratic has the two simple roots \(\beta=0\)
and \(\beta=2P_A(0)\), and the Jacobian matrix of
\eqref{eq:nodal-matching} is block triangular. For the nodal
parameters the system thus has exactly \(2S_p\) solutions, all simple.
By Step~4 the number of solutions counted with multiplicity is
\(2S_p=M_p\) for all \((g_2,g_3)\). This proves (a). Moreover the
nodal parameters do not belong to \(\mathcal N\). So \(\mathcal N\) is
a proper algebraic subset of the plane, and outside \(\mathcal N\)
every solution is simple. A nonzero polynomial cannot vanish on a
nonempty open subset of the plane, so \(\mathcal N\) does not contain
all the lattices. It contains only those with finitely many values of
\(j\), as we now show.

The equations \(E_i\) and their Jacobian determinant are
weighted-homogeneous, so \(\mathcal N\) is invariant under
\((g_2,g_3)\mapsto(\lambda^4g_2,\lambda^6g_3)\), \(\lambda\ne0\). On
lattices this is the rescaling \(\Lambda\mapsto\lambda^{-1}\Lambda\),
which maps pairs to pairs and simple pairs to simple pairs. If a
polynomial vanishes on \(\mathcal N\), then so does each of its
weighted-homogeneous components. Since \(\mathcal N\) is a proper
algebraic subset, it therefore lies in the zero set of a nonzero
polynomial that is a monomial in \(g_2,g_3\) times a product
\(\prod_i(g_2^3-c_ig_3^2)\). On the lattices this zero
set consists of finitely many values of \(j\): \(j=0\) on \(g_2=0\),
\(j=1728\) on \(g_3=0\), and \(j=1728c_i/(c_i-27)\) on
\(g_2^3=c_ig_3^2\). Two lattices with the same \(j\) differ by a
rescaling, so \(\mathcal N\) contains both or neither. Hence the
lattices in \(\mathcal N\) are exactly those whose invariant
\(j(\Lambda)\) lies in a finite set, which depends on \(p\).
This proves (b).
\end{proof}

\paragraph{Two backgrounds.}
The factor two in \(M_p\) retains the two choices of background in
\eqref{eq:shift}. That transformation preserves the evolution equation
and the lattice; it changes the operator \(P\) and the speed. For the
nodal parameters it exchanges the two roots of the quadratic in
\eqref{eq:nodal-matching}. Outside \(\mathcal N\) the two backgrounds of
every wave are distinct, because a solution with \(a_0=0\) is never
simple: Step~1 gives \(a_0+b\) as a polynomial in the \(u_j\) and
\(g_2,g_3\), and \(b\) enters the residual equations only through the
constant one, which gives \(a_0^2\) as another such polynomial; so the
Jacobian determinant contains the factor \(a_0\). On a generic lattice
the \(M_p\) pairs therefore represent \(M_p/2\) elliptic waves, up to
translations and Galilean shifts, each counted with both of its
backgrounds. On special lattices the two backgrounds may
coalesce, as \eqref{eq:kdvb-elliptic} shows for \(g_2=0\) and
Section~\ref{sec:ks} shows at order three.

The elliptic pairs of purely dispersive equations, \(2F_p\) of the
\(M_p\), are counted in Section~\ref{sec:reflection}
(Corollary~\ref{cor:elliptic-symmetric}), after the symmetry of their
profiles has been established; at order two they are both pairs of
\eqref{eq:kdvb-elliptic}, and \(M_2=2F_2=2\).
